% ============================================================================
% Appendix A11 — Fast--Slow Structure and the Quasi-Static Reduction
% ============================================================================
\section{Fast--Slow Structure and the Quasi-Static Reduction}
\label{app:fastslow}

Vitality appears to move faster than capital: the raw recovery rate is
$R=3$ against a depreciation rate $\mu=0.2$, and earlier drafts inferred a
``$\sim$15$\times$ timescale separation'' from the ratio of the naive
timescales $1/R=0.33$ and $1/\mu=5$. This appendix constructs the
quasi-static (slow-manifold) reduction formally, proves that its
equilibrium coincides \emph{exactly} with the full interior equilibrium,
and then \emph{measures}---rather than assumes---the quality of the
reduction. The honest finding: the separation at the equilibrium is a
factor $\approx3$, not $15$, and because the linearization there has a
complex eigenvalue pair, the fast and slow modes \emph{mix}---a strict
fast--slow decomposition fails precisely where the damped spiral lives.
The reduction is a legitimate intuition-builder for the slow capital
drift far from equilibrium; it is not a substitute for the
two-dimensional analysis, and none of the global results (in particular
the proof of Theorem~\ref{thm:gas} in Appendix~\ref{app:gas}) relies on
it. All quantities below are executable via \texttt{dlvt.fastslow} and
pinned by \texttt{tests/test\_fastslow.py}.

\subsection{The quasi-equilibrium and the slow manifold}

Freeze the capital at a level $C$, hence the complexity at
$O=O(C)=O_0+\beta C^{\eta}$ from \eqref{eq:complexity}. The vitality
equation \eqref{eq:dVdt} then relaxes to its quasi-equilibrium
$V_{\mathrm{qe}}(O)$, the unique positive root of the quadratic
\eqref{eq:vqe-quadratic},
\[
\frac{R}{\Vmax}V^{2}
-\Bigl(R-\frac{R\varepsilon}{\Vmax}-\delta O^{\gamma}\Bigr)V
-R\varepsilon=0 ,
\qquad
V_{\mathrm{qe}}(O)=\frac{-b+\sqrt{b^{2}+4\,(R/\Vmax)\,R\varepsilon}}{2R/\Vmax},
\]
with $b=-(R-R\varepsilon/\Vmax-\delta O^{\gamma})$; the product of the
roots is $-\varepsilon\Vmax<0$, so the ``$+$'' branch is always the
positive one, and $V_{\mathrm{qe}}$ is strictly decreasing along
$C\mapsto V_{\mathrm{qe}}(O(C))$ (Proposition~\ref{prop:nullcline}). Define the
\emph{slow manifold} as the graph
\begin{equation}
\label{eq:slowmanifold}
\mathcal{M}
\;=\;
\bigl\{(V,C): V=V_{\mathrm{qe}}(O(C)),\ C\ge0\bigr\}.
\end{equation}
By construction $\mathcal{M}$ is not an approximation to the
$V$-nullcline of the full system---it \emph{is} the $V$-nullcline, since
the quasi-equilibrium condition at frozen $O$ is literally $\dot V=0$.
Slaving $V$ to $\mathcal{M}$ in the capital equation
\eqref{eq:dCdt} yields the reduced one-dimensional slow flow
\begin{equation}
\label{eq:reducedflow}
\dot C
\;=\;
C\Bigl(\frac{\alpha\,V_{\mathrm{qe}}(O(C))}{1+\varphi O(C)}-\mu\Bigr)
\;=:\;
C\,g(C),
\end{equation}
the quasi-steady-state approximation familiar from enzyme kinetics and
ecological modeling \citep{murray2002mathematical}. Implementation:
\texttt{dlvt.fastslow.v\_quasi\_equilibrium},
\texttt{slow\_manifold}, \texttt{reduced\_slow\_rhs},
\texttt{simulate\_reduced}; the closed-form root satisfies the vitality
equation to residual $<10^{-13}$ across $O\in\{1,5,9,30\}$
(\texttt{test\_v\_quasi\_equilibrium\_satisfies\_vitality\_equation}).

\subsection{The reduced equilibrium is the full equilibrium---exactly}

\begin{proposition}[Zero equilibrium bias of the reduction]
\label{prop:reducedeq}
$C^{\dagger}>0$ is an equilibrium of the reduced flow
\eqref{eq:reducedflow} if and only if
$\bigl(V_{\mathrm{qe}}(O(C^{\dagger})),\,C^{\dagger}\bigr)$ is an interior
equilibrium of the full system
\eqref{eq:dVdt}--\eqref{eq:dCdt}. In particular the reduced flow has
exactly one positive equilibrium, $C^{\dagger}=C^{*}$, with
$V_{\mathrm{qe}}(O(C^{*}))=V^{*}$.
\end{proposition}

\begin{proof}
$g(C^{\dagger})=0$ is precisely the capital-nullcline condition
\eqref{eq:cnullcline} evaluated at $V=V_{\mathrm{qe}}(O(C^{\dagger}))$, and
membership in $\mathcal{M}$ is the $V$-nullcline condition. A point
satisfying both nullcline conditions is an interior equilibrium of the
full system, and conversely every interior equilibrium lies on both
nullclines, hence on $\mathcal{M}$ with $g=0$. Uniqueness is inherited
from Theorem~\ref{thm:uniqueness}.
\end{proof}

Unlike generic quasi-steady-state reductions, which displace equilibria
by $O(\epsilon)$ in the separation parameter, this reduction has
\emph{zero} equilibrium bias, because the manifold used for the slaving
is the exact nullcline rather than an asymptotic expansion of it.
Numerically, a Brent solve of $g(C)=0$ near $C^{*}$ reproduces
$(V^{*},C^{*})=(4.70253,\,32.0337)$ to a mismatch of
$7\times10^{-15}$
(\texttt{test\_reduced\_equilibrium\_coincides\_with\_full\_interior\_equilibrium}).

\subsection{How separated are the timescales, really?}

The raw-rate comparison ($1/R=0.33$ versus $1/\mu=5$, a factor $15$) is
the wrong calculation twice over.

\emph{First correction: linearize, don't compare raw coefficients.} Near
the equilibrium, the vitality relaxation rate is not $R$ but
\[
|J_{VV}|
=\frac{R}{\Vmax}
+\frac{\delta O^{*\gamma}\varepsilon}{(V^{*}+\varepsilon)^{2}}
=0.30704
\qquad(\text{Appendix~\ref{app:jacobian}}),
\]
because the logistic recovery term relaxes at rate $R/\Vmax$, not $R$.
Symmetrically, the capital rate is not $\mu$ but
\[
|J_{CC}|
=\frac{\alpha C^{*}V^{*}\varphi\,\beta\eta\,C^{*\,\eta-1}}
       {(1+\varphi O^{*})^{2}}
=0.10218 ,
\]
because on the capital nullcline the impact inflow cancels the
depreciation ($\alpha V^{*}/(1+\varphi O^{*})=\mu$ exactly), leaving only
the complexity-feedback term. The measured diagonal ratio is
\[
\biggl|\frac{J_{VV}}{J_{CC}}\biggr| \;=\; 3.005 ,
\]
a factor of $\sim3$, not $15$
(\texttt{test\_eigen\_separation\_is\_about\_3\_not\_15\_and\_modes\_mix}).

\emph{Second correction: a diagonal ratio of $3$ does not deliver modes.}
A fast--slow decomposition in the sense of singular perturbation theory
requires a real spectral gap between a fast and a slow eigenvalue. Here
the discriminant of the Jacobian is negative
($(\operatorname{tr}J)^{2}-4\det J=-0.4388$, Appendix~\ref{app:jacobian}),
so the eigenvalues are the complex pair
$\lambda_{1,2}=-0.205\pm0.331i$: both ``modes'' share one decay rate
($|\mathrm{Re}\,\lambda|=0.205$) and one rotation rate, and their common
modulus is $\sqrt{\det J}=0.389$. Near the equilibrium there are no
separate fast and slow directions at all---the coordinates mix into a
single damped-spiral mode \citep{strogatz2015nonlinear}. That is not a
technicality; it is \emph{why} the approach to $(V^{*},C^{*})$ is an
overshoot-and-return spiral rather than a monotone relaxation.

\subsection{Measured reduction error}

The consequence is quantitative and measurable. Integrating the full
system and the reduced flow \eqref{eq:reducedflow} from the same
$C_{0}$ (the reduced trajectory necessarily starts on $\mathcal{M}$)
and comparing $C(t)$ on a common grid over $[0,120]$
(\texttt{dlvt.fastslow.reduction\_error}) gives, at baseline:

\begin{center}
\begin{tabular}{lccc}
initial condition $(V_0, C_0)$ & max rel.\ error in $C(t)$ &
mean rel.\ error & equilibrium mismatch\\
\hline
$(8,\ 5)$    & $15.3\%$ & $1.3\%$ & $7\times10^{-15}$\\
$(2,\ 60)$   & $20.6\%$ & $1.2\%$ & ---\\
$(9.5,\ 80)$ & $20.4\%$ & $1.6\%$ & ---\\
\end{tabular}
\end{center}

The error is \emph{not} uniformly ``a few percent'': from $(8,5)$ it
peaks at $15.3\%$ at $t\approx9$, exactly during the overshoot---the full
trajectory overshoots to $C\approx36.5>C^{*}$ while the reduced flow,
being one-dimensional and therefore monotone, approaches $C^{*}$ from
below and can never overshoot. Away from the spiral the reduction is
excellent: the mean error is $\sim1\%$, and for $t>60$ the two solutions
agree to $<10^{-5}$ (both sit at the---identical---equilibrium). The
pinning test asserts the measured band honestly: maximum relative error
below $0.25$ \emph{and above} $0.05$, so that a future ``improvement''
cannot silently claim the reduction to be better than it is
(\texttt{test\_reduction\_tracks\_slow\_dynamics\_with\_measured\_error}).

\subsection{The basin portrait}

Theorem~\ref{thm:gas} makes a stark prediction: within the trapping
region the basin of the interior equilibrium is the \emph{entire} open
half-plane $\{C>0\}$---no basin boundary, no competing attractor, no
escape. The colored basin portrait makes this executable: on a
$40\times40$ grid of initial conditions
$(V_{0},C_{0})\in[0.2,\Vmax]\times[0.2,90]$, the full system is
integrated and the first time after which the trajectory enters and
\emph{remains} within the box $\{|V-V^{*}|<0.5,\ |C-C^{*}|<0.5\}$ is
recorded (\texttt{dlvt.fastslow.basin\_portrait\_grid}). All $1600$
initial conditions converge---zero exceptions---with times ranging from
$0$ (starting inside the band) to $30.0$, median $19.0$; the median is,
satisfyingly, one rotation period $2\pi/0.331\approx19.0$ of the spiral,
since most trajectories need one final loop to settle into the band. The
portrait is a smooth gradient of finite times, which is exactly what a
single global attractor looks like; a basin boundary would appear as a
discontinuity, and none exists
(\texttt{test\_basin\_grid\_all\_positive\_capital\_ics\_converge}).
The only slow region hugs the invariant axis $C\approx0$, where
trajectories linger near the axis saddle $(9.934,\,0)$ of
Appendix~\ref{app:regularization} before the unstable direction ejects
them into the interior.

\subsection{Remarks}

\begin{enumerate}
\item \emph{What the reduction is for.} The reduced flow
\eqref{eq:reducedflow} explains the slow architecture of the transient
in one line: capital drifts up (or down) the $V$-nullcline at the rate
$g(C)$, and the sign of $g$---positive below $C^{*}$, negative
above---is the entire one-dimensional story. This is the right mental
model for the long build-up phase of Figure~\ref{fig:temporal} and for
comparative statics along the manifold.
\item \emph{What it is not for.} No one-dimensional autonomous flow can
oscillate, so the reduction structurally cannot reproduce the
overshoot-and-return approach, the $\sim19$-unit rotation period, or any
quantity that depends on the spiral (peak overshoot, settling time into
a band). Those require the full planar system.
\item \emph{Independence of the global proof.} The proof of
Theorem~\ref{thm:gas} (Appendix~\ref{app:gas}) uses a trapping rectangle,
a Dulac certificate, and Poincar\'e--Bendixson; it makes no use of
timescale separation, and its validity is unaffected by the failure of
the strict fast--slow decomposition near the equilibrium. The
``$\sim$15$\times$ separation'' phrasing of earlier drafts is retracted
in favor of the measured statement: diagonal ratio $\approx3$, complex
pair, mode mixing.
\end{enumerate}
